\section{Matching the actual cloud data and denominator budgets}
\label{sec:cloud-parameters}

\status{The numerical data for the actual signed-chain test are now
constructed from the Type-II hypotheses. The floor in the slack is retained.
The outgoing realization used at the end is the separate OR argument;
ordinary source-symbol realization alone would not suffice.}

\paragraph{CP1: one set of budgets, with the correct floor.}
Write $k=m+1$, $N=|TT|$, $Q=p^{G'}$, and
\[
 L=Nk^2\log_p a,\qquad
 \ell=rNc_0+2rn'c'',\qquad c_P=c_*+\ell.
\]
Besides the actual orbit, weight and frame conditions, the hypotheses give
\[
\begin{gathered}
 r\ge2,\quad r\le Y,\quad N_1+3r\le k,\quad
 256k^2NY^2\le p^{L_1-1},\\
 CL\le D_p,\quad 8(D_p+1)\le c_*,\quad
 Ck^3r^2NQ^5c_PY\le a .
\end{gathered}                                                     \tag{CP.1}
\]
The last inequality, with $C\ge100$, gives
$a\ge\max(2,p,k,4N,r)$ and $L\ge1$; also $c_P>8CL$.
For $B>0$ put $u=B/(Yk^2)$ and $s(B)=\lfloor u\rfloor$.
If $B\ge a/Q^2$, then
\[
 u\ge Ckr^2NQ^3c_P\ge2,\qquad s(B)\ge u/2.
\]
The gap budget in (CP.1) gives
$p^{L_1+1}\ge48rk^2Y$, whence
\[
 3rp^{-L_1-1}B\le u/16\le s(B)/8.                     \tag{CP.2}
\]
If $B\ge a/Q$, then $u\ge Ckr^2NQ^4c_P\ge100r^2c_P$.
The last target is an integer, so
\[
                         s(B)\ge100r^2c_P.            \tag{CP.3}
\]
This proves directly the two numerical comparisons used downstream;
no rounding term is omitted.

\paragraph{CP2: the literal chain and outgoing margins.}
Let $z_j=f_{j,r-1}$, $D_j=z_j-f_{j,0}$,
$g_j=z_{j+1}-z_j$, and $P_j=p^{g_j}$.
For a cut in segment $j+1$, its oldest normalized degree $B$ satisfies
\[
 B\ge a/p^{D_j+g_j}\ge a/Q.
\]
The periodic wrap uses a common translation by $n'$ and leaves
$D_j+g_j$ unchanged. The stronger frame gap places every intermediate
letter of segment $j$ more than $L_2$ before $f_{j,0}$.
For its digit-block plane degree $d_x$, the exact degree identity gives
$p^{z_j-x}d_x\le2a$. Summing over distinct exponents and using $p\ge3$,
\[
 d_{LL,j}\le3ap^{-D_j-L_1-1},\qquad
 r d_{LL,j}/P_j\le s(B)/8.
\]
Since $26r\le25r^2$ for $r\ge2$, (CP.3) now proves
\[
 r d_{LL,j}/P_j+26rc_P\le3s(B)/8.                      \tag{CP.4}
\]
This is exactly the source's \texttt{SignedChainData.Hyp 26}.
For every non-old-frame cut the same oldest-degree bound gives
$s(B)\ge100r^2c_P\ge128c_P$. Together with the prescribed selected
degrees and zero degrees off $TT$, this is the complete
\texttt{MFH} condition required by DT10.

\paragraph{CP3: all logarithmic costs come from actual counts.}
The late letters in a single period inject into $TT$. Thus there are
at most $N$ true letters, at most $N$ chains, at most $N$ counted
links, and at most $2N$ normal projector cuts. Their block degrees are
at most $2a$. Put $h=\log_p a$. Since $a\ge p,4N,r,2$,
\[
 1+\log_p(d_x+r)\le4h,\qquad
 1+\log_p(2\sum_xd_x+r)\le4h.
\]
Consequently the literal definitions give
\[
 D_m\le4L,\quad D_l\le4L,\quad
 0\le k_p\le4L,\quad M_D:=D_m+D_l+3k_p\le20L.         \tag{CP.5}
\]
Also $\kappa_{\rm fl}=\sum_{t\in TT}k\operatorname{ilog}_p(d_t)
\le2Nkh\le20L$.
For each normal cut,
$2\sum_{x\in S}d_x+K-1\le4Na+k\le a^3$ and
$\binom K2\le k^2$. Hence its projector cost sums to $D_{\rm no}\le6L$.
Legendre's formula bounds the two multinomial deficits by
$3Nrh$ and $3NKh$. Their halves plus $D_{\rm no}$ are at most $9L$.
These are the actual \texttt{NumStep 20} and \texttt{SizeAt 9}
budgets. In particular $4M_D<c_P$.

\paragraph{CP4: the signed-chain object actually exists.}
Choose a unit coefficient of the primitive input polynomial and take
its digit class and plane/normal block degrees, as in DT8. Enumerate
the intermediate letters followed by the current frame. At a cut
write $B_u$ for the oldest degree, $s_u=s(B_u)$ and
\[
 T_u=\max\bigl(B_u/2,\ t_{1,u},\ t_{2,u}\bigr),
\]
where $t_{1,u}=a_{u+1}+m_*+s_u$ at an even letter and $B_u/2$
otherwise; $t_{2,u}=a_{u-1}+m_*+s_u$ after an odd letter and $B_u/2$
otherwise. Here $a_u$ is the first coordinate of the normal joint
highest weight. Use all old-frame carries and the first $i$
intermediate letters as chain $i$. Set its threshold to
$T_i-s_i/8$, its slack to $3s_i/8$, and each required link threshold
to $T_u-3s_u/8$.

These chains are strictly increasing. Link thresholds are positive:
$3s_u/8<B_u/2\le T_u$. To avoid any circular choice, first use zero
normal weights, select the joint projector by IW/DT, and then update
the thresholds to its actual weights. The projector depends on the
letter sets and degrees, not on these thresholds. DT supplies the
unit digit factor and its finite binomial approximations. The result
is literally \texttt{mkDt}, with every \texttt{Tied} identity satisfied.
IN/WT supply its incoming condition for the whole permitted word
span, since $C\ge100\ge16$. Thus the PG test applies to this datum.

\paragraph{CP5: the value bound is now unconditional on extra local budgets.}
With $f=\prod_td_t!$, DT's coefficient selection gives
$v_p(k_0)\le v_p(f)-n_0+9L$.
PG, with $(C_H,C_E,C_6,C_{18})=(26,2,24,60)$, gives a main pairing
of valuation at most $v_p(f)+29L$. The retained approximation error
has valuation at least $c_P-20L+v_p(f)-\kappa_{\rm fl}$.
Since $c_P>800L>69L$, its valuation is strictly higher than the main
term's. The normalized integral pairing therefore divides
\[
 p^{D_X},\qquad D_X=\lceil49L\rceil_+\le50L.           \tag{CP.6}
\]
The PG simplification $R=0$ is allowed, but does not remove the normal
digit-approximation error. DT10 now applies with the proved MFH
margins and decomposes the test into the actual permitted one-sided
symbols. This supplies \texttt{Proposition5At} and \texttt{NSStep}.

\paragraph{CP6: a uniform height, with constant $10$.}
By RK and the compatibility in EA, the initial symbols have normal
degree at most
\[
 X=10N\left(1088\min(N_1,2N)^2\log(a)\frac{a}{pk}
              +ap^{-L_1}+kQ\right).
\]
The fixed localizer multiplier has normal degree zero. Conjugating
an elementary operator by the chart substitution gives a linear
combination of operators which replace one variable by one variable;
each increases normal degree by at most one. The permitted word uses
at most $c_0$ operators at each of the $N$ positions. Hence
\[
 t_F=\lfloor X\rfloor_++Nc_0\le t_0(C,10)+Nc_0
\]
for $C\ge1088$. Every symbol component with normal degree $>t_F$
pairs to zero with the input, so discard it. The remaining components
have a common height bound independent of the primitive input.

\paragraph{CP7: realization and the normalized dual pairing.}
Apply OR to these one-sided symbols on the dual edge chart. Its
integral lift $F_o$ satisfies
$\Phi F_o=w^o\sigma+ w^{o+c_P}E$ and has order at least $o+c_P$
at every other plane. The opposite-edge Fock identification of IW
sends $w^{-o}F_o$ to an integral edge dual test after normalization by
$p^{\kappa_{\rm fl}}/f$. Legendre's inequality
$f\mid p^{\kappa_{\rm fl}}\alpha!$ makes its error integral.
Summing the tests gives exactly
\[
 \beta_{\rm fl}(\Theta F_1,T)=y+w^{c_P}z,\quad z\in O,
 \qquad yb=p^{D_X}.
\]
The lift has target height $\le t_F$ and other-edge order $c_P$.
Because $D_X\le D_p$ and $c_P\ge8(D_p+1)>2D_X$,
the value $y+w^{c_P}z$ still divides $p^{D_X}$: its product with $b$
is $p^{D_X}(1+w^{c_P-2D_X}bz)$, whose parenthesis is a unit.
This proves the actual \texttt{FlatAt} assertion, including its
integral error and every other-edge condition.

\paragraph{CP8: full cloud test and quantifier order.}
CB and the primitive-array argument CL3 apply to every coefficient
array having a unit entry. The normalized contraction of CL/IW has
cost $\kappa\le4L$ and shifts height and other-edge order by $\ell$.
If its scalar is $c$, the exact identity is
$c\beta(U,\Lambda T)=p^\kappa\beta_{\rm fl}(\Theta F_1,T)$.
The scalar need not be a unit. If the latter flat value times $b'$
equals $p^{D_X}$, then
\[
 \beta(U,\Lambda T)(cb')=p^{\kappa+D_X},\qquad
 D_1:=\kappa+D_X\le54L.                               \tag{CP.7}
\]
Other-edge order becomes $c_P-\ell=c_*$ and target height is at most
$t_0(C,10)+Nc_0+\ell$. Thus the exact cloud interface has
$(C_{t_0},C_{D_1})=(10,54)$ for every $C\ge4129$, since OR's actual
realization threshold is $100$. All constants are chosen before the local
field, edge, basis and coefficient array; the test itself may depend
on the array. The cloud result is a local theorem under the stated
Type-II hypotheses, not the later global supply of those hypotheses.
